\section{Predetermined mobility: the optimized moment threshold}
\label{sec:predetermined}
A predetermined design is selected before the offset $c$ is known. Its
optimal unrooted moment is
\begin{equation}\label{eq:blind-value}
 \mathcal P_q(M)=\inf_{D\in L^1(\Gamma),\ D\ge0,\ \int_\Gamma D=M}
       \frac14\int_{-2}^2J_c(D)^q\dd c,\qquad q>0.
\end{equation}
The admissible fields may depend on $M$ and may have arbitrarily fine
structure. In particular, none of the lower bounds below assumes a
smooth limiting design. The scalar convention is \eqref{eq:J-definition};
a degenerate coefficient is not automatically assigned a diffusion process.

\begin{theorem}[Sharp predetermined moments]\label{thm:blind-sharp}
Let $\alpha_q=(6q-4)/(q+4)$ and let $C_0$ be the harmonic coefficient
in \cref{prop:harmonic}. For each fixed $q>0$, as $M\downarrow0$,
\begin{equation}\label{eq:blind-three}
 \mathcal P_q(M)\sim
 \begin{cases}
 K_q M^{-q/4},&0<q<8/5,\\
 K_{\rm crit}M^{-2/5}[\log(1/M)]^{7/5},&q=8/5,\\
 2^{(11q-12)/7}\mathcal S_q M^{(2-3q)/7},&q>8/5,
 \end{cases}
\end{equation}
where
\begin{align}
 K_q&=2^{q-3}C_0^q
 \left[2\Bfun\!\left(\frac{1-\alpha_q}{2},\frac12\right)\right]^{1+q/4},
                                      \label{eq:blind-Kq}\\
 K_{\rm crit}&=\frac{(2C_0)^{8/5}}8\left(\frac47\right)^{7/5}.
                                      \label{eq:blind-Kcrit}
\end{align}
The finite positive constant $\mathcal S_q$ is the exact whole-line
variational value defined in \eqref{eq:Sq}. Its role in
\eqref{eq:blind-three} includes unrestricted lower and exact-budget
recovery bounds; it is not an inference from scaling alone.
For the corresponding full-bulk optimization with no observation,
each right-hand side is multiplied by $\chi^q$ under the hypotheses
of \cref{thm:optimal-transfer}.
\end{theorem}

Compared with \cref{thm:uniform-moments}, spatial design moves the
moment threshold from $4/3$ to $8/5$. The conclusion concerns a fixed
positive moment order. Taking the $q$th root divides the budget powers
by $q$; it does not change the minimizing fields. Orders below one
still give convex design problems, although they should not be described
as conventional risk aversion.

\subsection{Convexity and lower certificates without regularity assumptions}
\begin{lemma}[Convexity and symmetry]\label{lem:blind-convex}
For every $q>0$, $D\mapsto J_c(D)^q$ is convex in the extended sense.
The infimum in \eqref{eq:blind-value} is unchanged if designs are
required to be even and $\pi$-periodic.
\end{lemma}
\begin{proof}
The quotient formula \eqref{eq:J-quotient} gives
\[
 J_c(D)^q=\sup_f
 \frac{|\int f|^{2q}}{\bigl(\int D|f'|^2+\int k_cf^2\bigr)^q}.
\]
Every denominator is affine in $D$, and $x\mapsto x^{-q}$ is convex
on $(0,\infty)$, with its lower semicontinuous extension at zero.
The supremum is therefore convex. Reflection preserves $k_c$;
translation by $\pi$ replaces $c$ by $-c$, which has the same law.
Averaging over these four transformations preserves the budget and
cannot increase the expected moment.
\end{proof}

A useful lower bound follows by sliding a compact test along the possible
root positions. Fix a nonnegative $\psi\in C_c^\infty(-1,1)$ with
$\int\psi>0$. Near the fold at zero write $c=-\cos u$, and retain
$u\in[r,3r/2]$. Let $I_r$ be a fixed enlargement of this interval,
chosen inside $(r/2,2r)$, and put $m=\int_{I_r}D$.
For $0<m\le c_1r^7$, choose
$\ell=\delta_*(m/r^3)^{1/4}$, with fixed sufficiently small $\delta_*$.
The supports of $\psi((s-u)/\ell)$ lie in $I_r$, and
$k_c(s)\le Cr^2(s-u)^2$ there. Thus
\[
 J_c(D)^q\ge
 \frac{c_q\ell^{2q}}{[B(u)+Cr^2\ell^3]^q},\qquad
 B(u)=\ell^{-2}\int D(s)|\psi'((s-u)/\ell)|^2\dd s.
\]
Fubini gives $r^{-1}\int_r^{3r/2}B(u)\dd u\le Cm/(r\ell)$.
Jensen's inequality applies to the negative power for every $q>0$.
Since $\dd c=\sin u\dd u\asymp r\dd u$, it yields
\begin{equation}\label{eq:blind-shell}
 \E[J_c(D)^q;\ u\in[r,3r/2]]
       \ge c_q r^{2-5q/4}m^{-q/4}.
\end{equation}
If $m=0$, applying the same argument with widths tending to zero
makes the integrated lower bound infinite. This covers that case
without assuming an inverse exists for the degenerate field.

One fixed regular interval proves $\mathcal P_q(M)\ge c_qM^{-q/4}$.
At $q=8/5$, use geometrically separated intervals with disjoint spatial
enlargements, at radii between a fixed $r_*$ and $CM^{1/7}$.
There are $N\asymp\log(1/M)$ intervals; each satisfies the small-mass
condition because $m_j\le M$. Summing \eqref{eq:blind-shell} gives
\[
 \E[J_c(D)^{8/5}]\ge c\sum_{j=1}^Nm_j^{-2/5}
       \ge cN^{7/5}M^{-2/5}.
\]
Above the threshold, a fixed bump $\psi(s/R)$ at the fold, with
$R=M^{1/7}$, suffices. For $|c+1|\le cR^2$, its source is of order
$R$ and its derivative and reaction energies are at most
$CM/R^2$ and $CR^5$. Consequently
\begin{equation}\label{eq:blind-fold-bump}
 J_c(D)\ge cR^{-3},\qquad
 \E[J_c(D)^q]\ge c_qR^{2-3q}=c_qM^{(2-3q)/7}.
\end{equation}
These elementary certificates already identify the powers and the
critical logarithm from below. Sharper certificates are needed for
the coefficients.

\subsection{Globally admissible graded designs}
\begin{lemma}[Graded response bounds]\label{lem:graded-response}
Let $r(s)$ be distance to $\{0,\pi\}$, let $-1<\alpha<6$, and set
\begin{equation}\label{eq:graded-design}
 D(s)=a_R(r(s)+R)^{-\alpha},\qquad
 a_R\asymp R^{6+\alpha},\qquad 0<R\ll1.
\end{equation}
Constants in the comparison may depend on $\alpha$ and the two
constants in $a_R\asymp R^{6+\alpha}$. Near a fold put $t=1-|c|$.
Uniformly over $c\in[-2,2]$,
\begin{align}
 J_c(D)&\le Ca_R^{-1/4}r^{-3/2+\alpha/4},
       &&t>0,\quad r=\sqrt t\ge CR,\label{eq:graded-root}\\
 J_c(D)&\le CR^{-3},&&|t|\le CR^2,\label{eq:graded-core}\\
 J_c(D)&\le C|t|^{-3/2},&&t<0,\quad |t|\ge CR^2.
                                                   \label{eq:graded-out}
\end{align}
On compact simple-root parameter sets the bound is $Ca_R^{-1/4}$;
on compact rootless sets it is bounded. The same bounds hold with
$r(s)+R$ replaced by $|\sin s|+R$, or when the actual mobility is
bounded below by a fixed positive multiple of the displayed design.
\end{lemma}
\begin{proof}
For separated roots near a fold, take intervals of radius a fixed
small multiple of $r$. On them the mobility is comparable to
$a_Rr^{-\alpha}$, and the potential is bounded below by a constant
times $r^2$ times squared distance from the root. The natural-endpoint
harmonic estimate in \cref{lem:interval-response,prop:harmonic}, after
rescaling, gives the first bound for the root intervals. Their
oscillator-scaled lengths are bounded below, because
$a_R/r^{6+\alpha}\le C$; increasing the fixed threshold $r/R$
ensures this condition. The complement contributes at most $Cr^{-3}$
by its reciprocal potential. This is no larger than the root bound,
since their ratio is at most $C(R/r)^{(6+\alpha)/4}$.

In the central window use an interval of radius $C_2R$, large enough
to contain all roots for the stated offset window. Its mobility is
comparable to $R^6$. Rescaling $s=Rx$ gives derivative and potential
energies of order $R^5$ and source of order $R$, with a uniformly
anchored quartic interval problem. Its response is $O(R^{-3})$.
The complement has reciprocal-potential integral $O(R^{-3})$.
On the rootless side, discard the derivative term and integrate
$C(|t|+s^2)^{-2}$; the answer is $C|t|^{-3/2}$.
Fixed parameter sets are handled by the same harmonic localization
or by a positive potential bound. Comparability of $|\sin s|$ and
$r(s)$ proves the replacement statement; monotonicity in the
mobility proves the final assertion.
\end{proof}

Choose $a_R$ to impose the exact budget in \eqref{eq:graded-design}.
For $\alpha=\alpha_q$ use
\begin{equation}\label{eq:graded-cutoffs}
 R=\begin{cases}
 M^{1/(6+\alpha_q)},&q<8/5,\\
 [M/\log(1/M)]^{1/7},&q=8/5,\\
 M^{1/7},&q>8/5.
 \end{cases}
\end{equation}
Indeed, the normalization integral is comparable respectively to
$1$, $\log(1/R)$, and $R^{1-\alpha_q}$, so
$a_R\asymp R^{6+\alpha_q}$ in all three cases.
The root contribution to the moment is bounded by
\begin{equation}\label{eq:graded-integral}
 Ca_R^{-q/4}\int_R^{r_*}r^{1-3q/2+\alpha_q q/4}\dd r,
 \qquad
 2-\frac{3q}{2}+\frac{\alpha_q q}{4}
       =\frac{8-5q}{q+4}.
\end{equation}
It is bounded at its lower endpoint below $8/5$, logarithmic at
$8/5$, and controlled by that endpoint above $8/5$.
The core contributes $O(R^{2-3q})$; integrating
\eqref{eq:graded-out} gives a bounded term, a logarithm, or that same
power, all within the claimed orders. This proves matching upper
orders in \eqref{eq:blind-three}. For $q>8/5$, any
$1<\alpha<6-8/q$ gives the supercritical order.

\subsection{Sharp coefficients below and at the threshold}
We next retain the constants lost in \eqref{eq:blind-shell}.
Write $\beta=q/4$ and $W_q(r)=|\sin r|^{1-3q/2}$.
Let $E$ be the paired root arcs corresponding to a compact symmetric
interval of offsets inside $(-1,1)$, and define
\[
 Z_E=\int_E W_q^{1/(1+\beta)},\qquad
 d_E(r)=Z_E^{-1}W_q(r)^{1/(1+\beta)}.
\]
Extend this reference formula smoothly slightly beyond $E$. It
specifies test functions, not a constraint on the competing design.
Choose a real compact smooth $\psi$ with positive integral and put
\[
 I=\int\psi,\quad T=\int|\psi'|^2,\quad
 U=\int x^2\psi^2,\quad Q_\psi=T+U,\quad j_\psi=I^2/Q_\psi.
\]
For a root $r$ set $a(r)=\sin^2r$ and
\begin{align*}
 \ell_r&=(Md_E(r)/a(r))^{1/4},&
 h_r(s)&=(Md_E(r)a(r))^{-1/2}\psi((s-r)/\ell_r),\\
 z_c&=M^{-1/4}d_E(r)^{-1/4}a(r)^{-3/4}.
\end{align*}
The paired test $h_c$ has source $2Iz_c$. Its reaction energy is
$2Uz_c+o(z_c)$ uniformly on these offsets. Take the reference energy
$E_c^0=2Q_\psi z_c$, without requiring it to be the energy of any
single global reference field. The supporting line of $x^{-q}$ at
$E_c^0$ gives, for every competing $D$,
\begin{equation}\label{eq:sharp-tangent}
 J_c(D)^q\ge(2j_\psi)^qz_c^q
 \left[1+\frac{qT}{Q_\psi}
 -\frac{q}{2Q_\psi z_c}\int D|h_c'|^2+o(1)\right].
\end{equation}
The error is independent of $D$; a negative right-hand side causes
no difficulty for this supporting inequality.

For clarity, the moving-root calculation, including its probability
factor, is recorded explicitly. The marked-root density is
$\rho(r)=|\sin r|/4$. Summing over both roots counts each offset
twice, so
\begin{equation}\label{eq:sharp-reference-average}
 \E[(2j_\psi)^qz_c^q;\ c\text{ retained}]
 =K_{\psi,E}M^{-\beta},\qquad
 K_{\psi,E}=2^{q-3}j_\psi^qZ_E^{1+\beta}.
\end{equation}
The derivative penalty in \eqref{eq:sharp-tangent} has density
\[
 \mathcal K_M(s)=\E\!\left[
 \frac{q(2j_\psi)^q}{2Q_\psi}z_c^{q-1}|h_c'(s)|^2;
 c\text{ retained}\right].
\]
Changing from a root position $r$ to $y=(s-r)/\ell_r$ gives
\begin{equation}\label{eq:sharp-kernel}
 \|\mathcal K_M\|_\infty
 \le\left[\frac{qT}{Q_\psi}K_{\psi,E}+o(1)\right]M^{-1-\beta}.
\end{equation}
Here is the calculation behind this uniform bound. The squared
amplitude, squared derivative scale, and root Jacobian combine to
$M^{-5/4}d_E^{-5/4}a^{-3/4}$ before multiplying by
$z_c^{q-1}$. The resulting spatial factor is
\[
 M^{-1-\beta}\rho\,d_E^{-1-\beta}a^{-3q/4}
   =\tfrac14M^{-1-\beta}Z_E^{1+\beta}.
\]
The derivative profile integrates to $T$. All coefficients vary by
$o(1)$ across a support, uniformly on the fixed arcs. With the
prefactor in $\mathcal K_M$, this gives precisely
\eqref{eq:sharp-kernel}. At an arc endpoint only part of the
nonnegative kernel is retained, so the same upper bound holds even
there and just outside the arc. It can therefore be integrated
against an arbitrary concentrated $D$.

Averaging \eqref{eq:sharp-tangent}, using $\int D=M$, cancels the two
$qT/Q_\psi$ terms. Thus
\begin{equation}\label{eq:sharp-liminf}
 \liminf_{M\downarrow0}M^{q/4}\mathcal P_q(M)\ge K_{\psi,E}.
\end{equation}
The quotient supremum of the compact tests is $C_0$ by
\cref{prop:harmonic}. For $q<8/5$, expanding $E$ to all root arcs
gives $Z_E\to\int_0^{2\pi}|\sin r|^{-\alpha_q}\dd r$, which equals
the beta integral in \eqref{eq:blind-Kq} and is finite.

For the upper coefficient, take the exact-budget fields
\begin{equation}\label{eq:sharp-subcritical-design}
 D_M(s)=M\frac{(|\sin s|+R)^{-\alpha_q}}
                 {\int_0^{2\pi}(|\sin v|+R)^{-\alpha_q}\dd v},
 \qquad R=M^{1/(6+\alpha_q)}.
\end{equation}
For each fixed simple-root offset, the normalized shape converges
smoothly near both roots, and harmonic bracketing gives its local
coefficient $2C_0$. The graded estimates furnish the uniform envelope
\begin{equation}\label{eq:subcrit-envelope}
 M^{q/4}J_c(D_M)^q\le C|1-|c||^{-\gamma_q},\qquad
 \gamma_q=\frac{7q}{2(q+4)}<1
\end{equation}
near either fold, apart from the single fold point. On the root side
this is \eqref{eq:graded-root}. In the core, $M\asymp R^{6+\alpha_q}$
makes its left side at most $CR^{-2\gamma_q}$, which is bounded by
the displayed envelope. On the rootless side outside the core,
$M\le C|t|^{(6+\alpha_q)/2}$ and \eqref{eq:graded-out} give the same
bound. Elsewhere the normalized moment is uniformly bounded.
Dominated convergence now proves the upper coefficient
\eqref{eq:blind-Kq}, including $4/3\le q<8/5$ where the uniform-design
envelope would not suffice. Explicitly, with
$d_q=|\sin s|^{-\alpha_q}/Z_q$ and
$Z_q=\int_0^{2\pi}|\sin s|^{-\alpha_q}\dd s$, the limiting integral is
\[
 2^{q-3}C_0^q\int_0^{2\pi}W_q(s)d_q(s)^{-q/4}\dd s
       =2^{q-3}C_0^q Z_q^{1+q/4}=K_q.
\]

At $q_*=8/5$, let $0<b<1/7$, retain root distances at least $M^b$
from the folds, and use
\[
 Z_{E_M}=4b\log(1/M)+O(1),\qquad
 d_{E_M}(r)=\frac1{Z_{E_M}|\sin r|}.
\]
The same certificate is uniform on these moving arcs because
\begin{equation}\label{eq:critical-small-error}
 \sup_{r\in E_M}\frac{\ell_r}{\operatorname{dist}(r,\{0,\pi\})}
 \le C\left[\frac{M}{Z_{E_M}M^{7b}}\right]^{1/4}\longrightarrow0.
\end{equation}
The Taylor reaction error is controlled by this ratio. The logarithmic
derivatives of every coefficient in the moving-root kernel are
$O(1/r)$ near a fold, so its relative coefficient and Jacobian errors
are controlled by the same ratio. Endpoint truncation remains
nonnegative. Therefore the cancellation in \eqref{eq:sharp-liminf}
holds with relative errors $o(1)$ even though $Z_{E_M}$ diverges.
Take the small-budget limit first, then $j_\psi\to C_0$, then
$b\uparrow1/7$. This gives \eqref{eq:blind-Kcrit} from below.

For the matching upper bound put
\[
 L=\log(1/M),\quad R=(M/L)^{1/7},\quad
 Z_R=\int_0^{2\pi}\frac{\dd s}{|\sin s|+R},\quad
 a_R=M/Z_R,\quad D_M(s)=\frac{a_R}{|\sin s|+R}.
\]
There are four one-sided approaches to the two folds, hence
\begin{equation}\label{eq:critical-normalization}
 Z_R=4\log(1/R)+O(1)\sim\tfrac47 L,
 \qquad a_R\sim\tfrac74R^7.
\end{equation}
Retain root-to-fold distances $r\ge r_0=RL$. Around each root use
intervals of half-length $\eta|\sin r|$, $\eta=L^{-1/2}$.
The mobility and potential differ from their frozen quadratic values
by relative $O(\eta+R/r)$, uniformly. The oscillator width divided
by $r$ is at most $C(a_R/r_0^7)^{1/4}=O(L^{-7/4})$; thus the
rescaled intervals have lengths at least of order $L^{5/4}$.
The complement contributes at most $C/(\eta r^3)$, whose ratio to
the main contribution is $O(L^{-5/4})$. Both sides of the bracketing
therefore give the uniform equivalent
\[
 J_c(D_M)=[1+o(1)]2C_0a_R^{-1/4}(1-c^2)^{-5/8},
       \qquad |c|\le\cos r_0.
\]
Its critical moment integral is
\begin{equation}\label{eq:critical-main-integral}
 [1+o(1)]\frac{(2C_0)^{8/5}}8a_R^{-2/5}
           \int_{E(r_0)}\frac{\dd s}{|\sin s|}
 \sim\frac{(2C_0)^{8/5}}2a_R^{-2/5}\log(1/R).
\end{equation}
The omitted root annuli $CR<r<RL$ cost
$O(a_R^{-2/5}\log L)$ by \eqref{eq:graded-root}; the core and the
rootless sides cost $O(a_R^{-2/5})$ by
\eqref{eq:graded-core}--\eqref{eq:graded-out}. They are negligible.
Substitution of \eqref{eq:critical-normalization} into
\eqref{eq:critical-main-integral} proves \eqref{eq:blind-Kcrit}.

\subsection{The supercritical whole-line problem and measure limits}
For a finite nonnegative measure $\nu$ on $\R$, define the scalar test
response and its integrated moment by
\begin{align}
 \mathscr J_\mu(\nu)&=\sup_{v\in C_c^\infty(\R)}
 \left\{2\int v\dd x-\int|v'|^2\dd\nu
                   -\int(x^2-\mu)^2v^2\dd x\right\},
                                          \label{eq:measure-response}\\
 \mathscr I_q(\nu)&=\int_\R\mathscr J_\mu(\nu)^q\dd\mu,\qquad
 \mathcal S_q=\inf_{d\ge0,\ d\in L^1(\R),\ \int d=1}
                         \mathscr I_q(d\dd x).
                                          \label{eq:Sq}
\end{align}
Measures are introduced only for compactness of design sequences.
Equation \eqref{eq:measure-response} does not specify a stochastic
process with measure-valued mobility. Measure relaxation and
semicontinuity of compact-test variational values are established
methods in reinforcement optimization; see, for example,
\citet[Proposition~4.1]{ButtazzoOudetVelichkov2015}. The arguments below
also address the noncompact domain, singular mass, and the integrated
fold response needed here.

\begin{lemma}[Singular mass, weak limits, and scaling]
\label{lem:measure-relaxation}
Write $\nu=d\dd x+\nu_s$ for its Lebesgue decomposition. Then
\begin{equation}\label{eq:singular-useless}
 \mathscr J_\mu(\nu)=\mathscr J_\mu(d\dd x),\qquad
 \mathscr I_q(\nu)=\mathscr I_q(d\dd x).
\end{equation}
The response is lower semicontinuous under vague convergence of
finite measures, pointwise in $\mu$; the integrated moment is also
lower semicontinuous. For $q>8/5$, $0<\mathcal S_q<\infty$ and
\begin{equation}\label{eq:fold-mass-scaling}
 \inf_{\int d=m}\mathscr I_q(d\dd x)
       =m^{-\beta_q}\mathcal S_q,\qquad
 \beta_q=\frac{3q-2}{7}>0,\quad m>0.
\end{equation}
The value at $m=0$ is infinite.
\end{lemma}
\begin{proof}
The inequality $\mathscr J_\mu(\nu)\le\mathscr J_\mu(d\dd x)$ is
immediate. For the reverse inequality fix a compact smooth $v$.
Inside a common compact interval take compact sets $K_n$ supporting
all but $o(1)$ of the singular mass relevant to $v'$. Choose open
neighborhoods $U_n$ with $|U_n|\to0$ and $\int_{U_n}d\to0$.
There are smooth $0\le\rho_n\le1$, zero near $K_n$ and one outside
$U_n$. Fix a smooth compact $\zeta$ of integral one and set
\[
 a_n=\int\rho_nv',\qquad v_n'=\rho_nv'-a_n\zeta,
 \qquad v_n(x)=\int_{-\infty}^xv_n'(y)\dd y.
\]
Their supports lie in a common compact set, $a_n=O(|U_n|)$,
and $v_n\to v$ uniformly. The absolutely continuous derivative
energies converge to $\int d|v'|^2$, while the singular derivative
energies tend to zero: the uncorrected derivative vanishes on $K_n$,
the remaining singular mass tends to zero, and the correction has
size $a_n$. The source and reaction integrals converge uniformly
on the common compact support. This proves \eqref{eq:singular-useless}
by taking the test supremum.

For a fixed test the derivative integral is continuous under vague
convergence. The supremum is therefore lower semicontinuous. The
response is Borel in $\mu$ (it is a supremum of continuous functions),
so Fatou proves the integrated assertion. A fixed nonnegative bump
on $[-1,1]$, optimized in amplitude, gives a positive response bound
for $|\mu|\le1$ uniformly when $\nu(\R)\le1$. Thus
$\mathcal S_q>0$.

For finiteness choose
\begin{equation}\label{eq:tail-choice}
 1<\alpha<\min\{2,6-8/q\},\qquad
 \rho_\alpha(x)=c_\alpha(1+|x|)^{-\alpha},\quad\int\rho_\alpha=1.
\end{equation}
The variable-mobility version of the bracketing proof in
\cref{lem:graded-response} gives
\begin{equation}\label{eq:line-tail-envelope}
 \mathscr J_\mu(\rho_\alpha)\le C
 \begin{cases}
 (1+\mu)^{-(6-\alpha)/8},&\mu\ge0,\\
 (1+|\mu|)^{-3/2},&\mu<0.
 \end{cases}
\end{equation}
For completeness, at large positive $\mu$ take intervals of radius
$c\sqrt\mu$ around $\pm\sqrt\mu$. There the mobility is comparable
to $\mu^{-\alpha/2}$ and the quadratic coefficient is comparable
to $\mu$, giving the displayed harmonic power. The complementary
reciprocal potential is $O(\mu^{-3/2})$. For bounded $\mu$, a fixed
core has a positive mobility lower bound and an anchored energy;
for negative $\mu$ use the reciprocal potential directly outside
a fixed compact set of parameters. Both powers in
\eqref{eq:line-tail-envelope} have integrable $q$th powers by
\eqref{eq:tail-choice} and $q>8/5$.

Finally, if $d_m(x)=m^{6/7}d(x/m^{1/7})$, then
$\int d_m=m\int d$ and a change of variables in the test supremum gives
\[
 \mathscr J_\mu(d_m\dd x)
   =m^{-3/7}\mathscr J_{\mu/m^{2/7}}(d\dd x).
\]
Integration in $\mu$ proves \eqref{eq:fold-mass-scaling} in both
directions. At zero mass the response is infinite for every $\mu>0$,
by bumps at a root of $(x^2-\mu)^2$ with widths tending to zero.
\end{proof}

In particular, relaxation to finite measures of mass at most one does
not reduce $\mathcal S_q$. Singular mass has no benefit for this
scalar convention, and loss of mass at infinity increases the
optimal value through \eqref{eq:fold-mass-scaling}. The statement
also implies existence of at least one unit-mass density minimizing
\eqref{eq:Sq}: take a minimizing sequence of measures, extract a
vaguely convergent subsequence, use lower semicontinuity and
\eqref{eq:singular-useless}, and then use the strict decrease of
$m^{-\beta_q}$ to rule out lost or singular mass. No uniqueness,
regularity, or explicit formula for that minimizer is asserted.

\begin{lemma}[Natural endpoints with a graded positive background]
\label{lem:graded-neumann-limit}
Suppose $d\in L^1(\R)$ and
$d(x)\ge a(1+|x|)^{-\alpha}$, where $a>0$ and $1<\alpha<2$.
For $\mu$ in a compact set, the whole-line response is finite.
On intervals $(-T_n,T_n)$ with $T_n\to\infty$, the free-endpoint
response converges to $\mathscr J_\mu(d\dd x)$. The conclusion holds
along converging $\mu_n$ and when the source and potential are
multiplied by weights $w_n$ that are uniformly bounded above and
below and converge locally uniformly to one; the derivative energy
remains $\int d|v'|^2$.
\end{lemma}
\begin{proof}
On a fixed core the lower bound on $d$ and a positive-potential
anchor give an $H^1$ bound from the energy. Outside that core the
potential is at least $cx^4$, uniformly in the parameters. Hence
$|\int w_nv|\le C\mathfrak e_n[v]^{1/2}$ and
\[
 \int_{L<|x|<T_n}|w_nv|\dd x
      \le CL^{-3/2}\mathfrak e_n[v]^{1/2}.
\]
This gives uniformly bounded maximizing energies and responses.
The finite-interval energy space consists of absolutely continuous
functions with finite weighted derivative energy. Smooth functions
are dense there: approximate the derivative in $L^2(d\dd x)$ by
smooth functions, correct its integral with a fixed smooth bump
when zero traces are required, and integrate. On compact intervals
the positive lower bound on $d$ makes derivative convergence imply
uniform convergence of the anchored functions. Thus the reaction
and source terms converge as well.

A sequence of maximizing functions has a locally weak $H^1$ limit
$v$ and locally weak weighted derivatives. Distributional derivatives
identify these limits, and lower semicontinuity bounds the limiting
energy. The source tails above make the source integrals converge.
To use $v$ in the whole-line smooth-test supremum, its membership in
that energy completion must also be checked. For $|x|$ large, the
one-dimensional Sobolev estimate on $[x-1,x+1]$ gives
\[
 |v(x)|^2\le C\left[|x|^{-4}+|x|^{\alpha/2-2}\right]
  \int_\R\{d|v'|^2+(y^2-\mu)^2v^2\}\dd y.
\]
Indeed the local $L^2$ norm is bounded by $C|x|^{-4}$ times the
energy, and the local derivative norm by $C|x|^\alpha$ times the
energy. Their product gives the second term. In particular $v$
is bounded. Cutting off at $R$ therefore adds derivative error at
most $CR^{-2}\|v\|_\infty^2\int_{R<|x|<2R}d$, tending to zero;
the original energy tails also vanish. The compact derivative
approximation just described then gives compact smooth convergence
in energy. The limiting variational value is at most the whole-line
response. Conversely, every compact smooth test is eventually an
admissible interval test and its weighted integrals converge.
This proves both inequalities and, by subsequences, the assertions
for varying parameters and weights.
\end{proof}

\subsection{Sharp supercritical localization and exact-budget recovery}
Put $b_0=1/2$, the coefficient in
$1-\cos s=b_0s^2+O(s^4)$, and first use the common scale
$\ell=(M/b_0^2)^{1/7}$. Given any sequence of designs of budget $M$,
first pass to a subsequence realizing the lower limit of its scaled
objective. If that limit is infinite there is nothing to prove.
Otherwise restrict this subsequence to the two cells
$|s-s_j|<\pi/2$, $s_j=0,\pi$, and
push the mobility measures to $x=(s-s_j)/\ell$, dividing their
masses by $M$. Along a further subsequence these measures converge vaguely
to $\nu_0,\nu_1$, with total masses $\theta_0+\theta_1\le1$.
This compactness follows by weak compactness on each compact interval
and a diagonal subsequence; total mass is bounded by one.

For offsets $c=-1+b_0\ell^2\mu$ near the first fold, or
$c=1-b_0\ell^2\mu$ near the second, take compact tests
$h(s)=b_0^{-2}\ell^{-4}v((s-s_j)/\ell)$.
Their rescaled potential tends locally uniformly to $(x^2-\mu)^2$.
The scaled source and derivative integrals therefore show
\begin{equation}\label{eq:fold-liminf-local}
 \liminf b_0^2\ell^3J_c(D_M)\ge\mathscr J_\mu(\nu_j).
\end{equation}
One may first use any one compact test and then take its supremum;
no exchange of an optimizer with the limit is needed. The two
parameter windows are disjoint for every fixed bounded $\mu$ range
once $M$ is small. Fatou on those windows, followed by expansion of
the range to $\R$, gives
\begin{align}
 \liminf M^{\beta_q}\E[J_c(D_M)^q]
 &\ge \tfrac14 b_0^{(3-8q)/7}
       [\mathscr I_q(\nu_0)+\mathscr I_q(\nu_1)]\notag\\
 &\ge \tfrac14 b_0^{(3-8q)/7}\mathcal S_q
       [\theta_0^{-\beta_q}+\theta_1^{-\beta_q}]\notag\\
 &\ge 2^{(11q-12)/7}\mathcal S_q.
                                                   \label{eq:fold-sharp-lower}
\end{align}
If a limiting measure has a singular part, its absolutely continuous
mass is no greater than $\theta_j$, so the middle inequality still
follows from \cref{lem:measure-relaxation}. A zero mass is interpreted
as infinite cost. The final minimization uses $\theta_0+\theta_1\le1$
and the strict convexity and decrease of $\theta^{-\beta_q}$: the
best allocation is $\theta_0=\theta_1=1/2$. This proves the lower
bound for all competitors, including sequences with concentration,
oscillation, or mass escaping the rescaled fold neighborhoods.

Here is a matching recovery construction for the same unrestricted
value $\mathcal S_q$. Take a unit-mass $d$ with finite cost arbitrarily
close to $\mathcal S_q$, and add $\varepsilon\rho_\alpha$ with
\eqref{eq:tail-choice}. Monotonicity cannot increase the cost.
Rescale spatially by the mass transformation in
\eqref{eq:fold-mass-scaling} to restore unit mass. The resulting
profile, denoted $\widetilde d$, satisfies
\begin{equation}\label{eq:near-optimal-positive-profile}
 \int\widetilde d=1,\qquad
 \widetilde d(x)\ge a_\varepsilon(1+|x|)^{-\alpha},\qquad
 \mathscr I_q(\widetilde d\dd x)
       \le(1+\varepsilon)^{\beta_q}\mathscr I_q(d\dd x).
\end{equation}
It may be averaged with its reflection, preserving these properties
and decreasing the cost by convexity. No regularity or existence of
a minimizing profile is needed for this construction.

Now give each fold mass $m=M/2$, put
$\ell=(m/b_0^2)^{1/7}$, and on its cell use the exact coordinate
\[
 x=\frac{2\sin((s-s_j)/2)}\ell,\quad
 T_\ell=\frac{\sqrt2}\ell,\quad
 w_\ell(x)=(1-\ell^2x^2/4)^{-1/2},\quad
 \dd s=\ell w_\ell(x)\dd x.
\]
The metric obeys $1\le w_\ell\le\sqrt2$ on the entire cell and
converges locally uniformly to one. Define the preliminary field
\begin{equation}\label{eq:fold-recovery-design}
 \widehat D_M(s)=b_0^2\ell^6 w_\ell(x)\widetilde d(x)
 \quad\hbox{on each cell}.
\end{equation}
Its mass in each cell is
$m\int_{-T_\ell}^{T_\ell}w_\ell^2\widetilde d\dd x=m[1+o(1)]$
by dominated convergence. Multiply the whole field by
$M/\int\widehat D_M=1+o(1)$ to obtain an exactly admissible $D_M$.
This normalization changes each response by a factor between
$1$ and $1+o(1)$ in whichever direction it acts: for a multiplier
$\lambda$, the full quadratic energy is bounded between
$\min(1,\lambda)$ and $\max(1,\lambda)$ times its old value.

Near the first fold, $c=-1+b_0\ell^2\mu$ gives the exact potential
$b_0^2\ell^4(x^2-\mu)^2$. The scaled derivative energy from
\eqref{eq:fold-recovery-design} is precisely
$\int\widetilde d|v'|^2$; the source and reaction acquire the weight
$w_\ell$. Bracketing into the two cells and applying
\cref{lem:graded-neumann-limit} gives, for every fixed $\mu$,
\begin{equation}\label{eq:fold-recovery-pointwise}
 b_0^2\ell^3J_{-1+b_0\ell^2\mu}(D_M)
       \longrightarrow\mathscr J_\mu(\widetilde d\dd x).
\end{equation}
The other cell has a potential bounded below for these offsets and
contributes $O(1)$ before rescaling. Compact supported tests give
the reverse bracket. The same argument holds at the second fold.

Pointwise localization is not by itself sufficient for the moment.
The lower background in \eqref{eq:near-optimal-positive-profile}
implies on the whole circle
\[
 D_M(s)\ge c_\varepsilon\ell^{6+\alpha}(r(s)+\ell)^{-\alpha}.
\]
Apply \cref{lem:graded-response}, using each fold for its parameter
half $c\le0$ or $c\ge0$. With $\mu$ the corresponding signed scaled
offset, it yields the uniform envelope
\begin{equation}\label{eq:recovery-global-envelope}
 b_0^2\ell^3J_c(D_M)\le C_\varepsilon
 \begin{cases}
 (1+\mu)^{-(6-\alpha)/8},&\mu\ge0,\\
 (1+|\mu|)^{-3/2},&\mu<0.
 \end{cases}
\end{equation}
For offsets bounded away from a fold but with roots, the first bound
is the regular-root bound: $\mu\asymp\ell^{-2}$, so both scales are
$\ell^{(6-\alpha)/4}$. For fixed rootless offsets the second bound
is the reciprocal-potential bound. Thus the envelope covers the
entire half of the offset range, including its outer end; extend
its left side by zero outside that range. Its $q$th power is
integrable by \eqref{eq:tail-choice}. Dominated convergence in
\eqref{eq:fold-recovery-pointwise} proves
\[
 \lim_{M\downarrow0}M^{\beta_q}\E[J_c(D_M)^q]
       =2^{(11q-12)/7}\mathscr I_q(\widetilde d\dd x).
\]
Finally choose $d$ arbitrarily close to the infimum and then let
$\varepsilon\downarrow0$ in \eqref{eq:near-optimal-positive-profile}.
Together with \eqref{eq:fold-sharp-lower}, this proves the
supercritical equivalent in \cref{thm:blind-sharp}.

\subsection{Physical transfer and the limits of the design statements}
All the moment values in \eqref{eq:blind-three} diverge faster than
a fixed power of $\log(1/M)$. The same-budget theorem
\cref{thm:optimal-transfer} therefore transfers their exact equivalents
to the physical full-bulk optimum, with coefficient $\chi^q$.
The graded trials have a positive background for each fixed $M$.
The whole-line recovery profiles also have a positive background
after transplantation to the compact wall. Alternatively, the
positive-background mixture in \eqref{eq:mixture} supplies admissible
physical trials without changing any leading value. These statements
hold at fixed positive bulk diffusivity and nonzero mean speed;
no additional joint vanishing-bulk limit is asserted.

The asymptotic designs are not certified finite-budget optimizers.
Nor does achieving the correct order force every design to concentrate
all mobility near the folds: a successful graded trial with budget
$M/2$, plus a uniform field of budget $M/2$, still achieves the
supercritical order by monotonicity. The sharp supercritical lower
bound instead identifies the mass allocation of asymptotically
optimal sequences at the level of the value: equality requires full
mass at the two rescaled folds and equal limiting mass at them.
It does not prove uniqueness or a closed form for the local minimizing
density. A prescribed mobility cap, fixed minimum background, or
minimum fabrication length changes the admissible small-budget
problem and is outside \cref{thm:blind-sharp}.
